\begin{tikzpicture}[
    node distance=0.28in and 0.62in,
    function/.style={
        draw=black!60,
        fill=black!3,
        rounded corners=2pt,
        minimum width=1.90in,
        minimum height=0.30in,
        text width=1.76in,
        align=center,
        font=\ttfamily\scriptsize
    },
    entry/.style={
        function,
        fill=green!8,
        draw=green!40!black,
        line width=0.65pt
    },
    lifecycle/.style={
        function,
        fill=purple!7,
        draw=purple!50!black
    },
    wall/.style={
        function,
        fill=blue!7,
        draw=blue!45!black
    },
    reused/.style={
        draw=orange!55!black,
        fill=orange!10,
        rounded corners=2pt,
        minimum width=1.16in,
        minimum height=0.22in,
        text width=1.04in,
        inner xsep=2pt,
        align=center,
        font=\ttfamily\tiny
    },
    flow/.style={-{Stealth[length=5pt]}, draw=black!65, line width=0.45pt},
    reusecall/.style={-{Stealth[length=4pt]}, draw=orange!60!black, dashed, line width=0.4pt},
    loop/.style={flow, rounded corners=0pt}
]

\providecommand{\ForceFlowNode}[2]{\hyperref[force:#1]{\strut #2}}

\node[entry] (main) {main(SourceID)\\{\scriptsize WALL\_MODEL\_PIPE}};
\node[lifecycle, below=of main] (start-reservoir)
    {StartReservoir(SourceID)};
\node[function, below=of start-reservoir] (start-state)
    {\ForceFlowNode{InitializeStartingContactState}{InitStartingContactState}};

\node[function, below left=0.34in and 0.74in of start-state] (particle-scan)
    {target scan};
\node[function, below=of particle-scan] (particle-process)
    {ProcessParticleCollision};
\node[function, below=of particle-process] (geometry)
    {\ForceFlowNode{GetParticleGeometry}{GetParticleGeometry}};
\node[function, below=of geometry] (particle-force)
    {AccumulateContactForce};

\node[wall, below right=0.34in and 0.74in of start-state] (wall-scan)
    {manual wall scan\\{\scriptsize wall flags 3--4}};
\node[wall, below=of wall-scan] (wall-process)
    {ProcessWallCollision};
\node[wall, below=of wall-process] (wall-geometry)
    {GetWallGhostGeometry};
\node[wall, below=of wall-geometry] (wall-force)
    {AccumulateContactForce};

\node[entry, below=0.55in of main |- particle-force.south] (velocity)
    {CalcVelocity(SourceID, totalForce)};
\node[entry, below=of velocity] (position)
    {CalcPosition(SourceID)};
\node[lifecycle, below=of position] (retire)
    {RetirePastXMax(SourceID)};

\node[reused, left=0.16in of geometry] (particle-position)
    {\ForceFlowNode{GetParticlePosition}{GetParticlePosition}};
\node[reused, left=0.16in of particle-force] (particle-stiffness)
    {GetContactStiffness};
\node[reused, right=0.16in of wall-geometry] (wall-position)
    {\ForceFlowNode{GetParticlePosition}{GetParticlePosition}};
\node[reused, right=0.16in of wall-force] (wall-stiffness)
    {GetContactStiffness};
\node[reused, left=of velocity] (velocity-mass)
    {\ForceFlowNode{GetParticleMass}{GetParticleMass}};
\node[reused, right=of velocity] (velocity-start)
    {\ForceFlowNode{GetStartFrameVelocity}{GetStartFrameVelocity}};

\draw[flow] (main) -- (start-reservoir);
\draw[flow] (start-reservoir) -- node[right,font=\scriptsize]{active source} (start-state);
\draw[flow] (start-state) -- (particle-scan);
\draw[flow] (start-state) -- (wall-scan);

\draw[flow] (particle-scan) -- node[right,font=\scriptsize]{mobile target} (particle-process);
\draw[flow] (particle-process) -- (geometry);
\draw[flow] (geometry) -- (particle-force);
\draw[loop] (particle-force.west) -- ++(-0.18in,0)
    |- node[pos=0.12,left,font=\scriptsize]{next target} (particle-scan.west);

\draw[flow] (wall-scan) -- node[right,font=\scriptsize]{next wall} (wall-process);
\draw[flow] (wall-process) -- (wall-geometry);
\draw[flow] (wall-geometry) -- (wall-force);
\draw[loop] (wall-force.east) -- ++(0.18in,0)
    |- node[pos=0.12,right,font=\scriptsize]{wall 3--4} (wall-scan.east);

\draw[flow] (particle-force.south) -- ++(0,-0.18in) -| (velocity.north);
\draw[flow] (wall-force.south) -- ++(0,-0.18in) -| (velocity.north);
\draw[flow] (velocity) -- (position);
\draw[flow] (position) -- (retire);

\draw[reusecall] (geometry) -- (particle-position);
\draw[reusecall] (particle-force) -- (particle-stiffness);
\draw[reusecall] (wall-geometry) -- (wall-position);
\draw[reusecall] (wall-force) -- (wall-stiffness);
\draw[reusecall] (velocity) -- (velocity-mass);
\draw[reusecall] (velocity) -- (velocity-start);

\node[font=\sffamily\scriptsize, text=black!65, above=0.08in of particle-scan]
    {particle contacts};
\node[font=\sffamily\scriptsize, text=blue!55!black, above=0.08in of wall-scan]
    {pipe walls};

\end{tikzpicture}
